\subsection*{分级提示}
\begingroup\small
\begin{enumerate}[leftmargin=*]
  \item 分开有限样本中心、概率收缩、缩放后极限形状与同类估计量的渐近方差。
  \item 大小在零假设下定义，功效在备择下定义；区间覆盖是重复抽样频率。
  \item 先画误差相关矩阵；最后问系数目标本身是否已被正确识别。
  \item 加减 $\mathbb E\hat\theta$，交叉项因中心化消失；收缩会同时改变偏差与方差。
  \item 对 $b$ 求梯度；把一阶条件改写为 $X^T(y-X\hat\beta)=0$。
  \item 将 $z$ 对 $x$ 的线性投影系数代入真实模型；标准误不改变概率极限。
  \item 分别按固定第一项和最大项构造区间，统计覆盖固定真值的比例。
  \item 先求 $e=y-X\hat\beta$，再计算 $X^T\operatorname{diag}(e_i^2)X$ 与两侧逆矩阵。
  \item 标准误与 $1/\sqrt n$ 成正比；分清总体标准差与均值标准误。
  \item 十次独立查看的误拒概率是 $1-0.95^{10}$；确认集不得参与选择。
  \item 检查遗漏变量、碰撞点、处理后控制、未来信息、函数形式与测量误差。
  \item 把检验族、选择、样本外、依赖标准误、效应、成本、容量和限制串成不可回流的流程。
\end{enumerate}
\endgroup
